\documentclass[10pt,a4paper]{amsart}
\usepackage{amsmath,amssymb,amsthm,graphicx}
\usepackage{verbatim}
\newenvironment{reference}{\comment}{\endcomment}
%\newenvironment{reference}{}{}
\newenvironment{slogan}{\comment}{\endcomment}
\newenvironment{history}{\comment}{\endcomment}

% For commutative diagrams we use Xy-pic
\usepackage[all]{xy}

% We use 2cell for 2-commutative diagrams.
\xyoption{2cell}
% Theorem environments.
%
\theoremstyle{plain}
\newtheorem{theorem}[subsection]{Theorem}
\newtheorem{proposition}[subsection]{Proposition}
\newtheorem{lemma}[subsection]{Lemma}

\theoremstyle{definition}
\newtheorem{definition}[subsection]{Definition}
\newtheorem{example}[subsection]{Example}
\newtheorem{exercise}[subsection]{Exercise}
\newtheorem{situation}[subsection]{Situation}

\theoremstyle{remark}
\newtheorem{remark}[subsection]{Remark}
\newtheorem{remarks}[subsection]{Remarks}

\numberwithin{equation}{subsection}

% Macros
%
\def\lim{\mathop{\mathrm{lim}}\nolimits}
\def\colim{\mathop{\mathrm{colim}}\nolimits}
\def\Spec{\mathop{\mathrm{Spec}}}
\def\Hom{\mathop{\mathrm{Hom}}\nolimits}
\def\Ext{\mathop{\mathrm{Ext}}\nolimits}
\def\SheafHom{\mathop{\mathcal{H}\!\mathit{om}}\nolimits}
\def\SheafExt{\mathop{\mathcal{E}\!\mathit{xt}}\nolimits}
\def\Sch{\mathit{Sch}}
\def\Mor{\mathop{\mathrm{Mor}}\nolimits}
\def\Ob{\mathop{\mathrm{Ob}}\nolimits}
\def\Sh{\mathop{\mathit{Sh}}\nolimits}
\def\NL{\mathop{N\!L}\nolimits}
\def\CH{\mathop{\mathrm{CH}}\nolimits}
\def\proetale{{pro\text{-}\acute{e}tale}}
\def\etale{{\acute{e}tale}}
\def\QCoh{\mathit{QCoh}}
\def\Ker{\mathop{\mathrm{Ker}}}
\def\Im{\mathop{\mathrm{Im}}}
\def\Coker{\mathop{\mathrm{Coker}}}
\def\Coim{\mathop{\mathrm{Coim}}}

% Boxtimes
%
\DeclareMathSymbol{\boxtimes}{\mathbin}{AMSa}{"02}

%
% Macros for moduli stacks/spaces
%
\def\QCohstack{\mathcal{QC}\!\mathit{oh}}
\def\Cohstack{\mathcal{C}\!\mathit{oh}}
\def\Spacesstack{\mathcal{S}\!\mathit{paces}}
\def\Quotfunctor{\mathrm{Quot}}
\def\Hilbfunctor{\mathrm{Hilb}}
\def\Curvesstack{\mathcal{C}\!\mathit{urves}}
\def\Polarizedstack{\mathcal{P}\!\mathit{olarized}}
\def\Complexesstack{\mathcal{C}\!\mathit{omplexes}}
% \Pic is the operator that assigns to X its picard group, usage \Pic(X)
% \Picardstack_{X/B} denotes the Picard stack of X over B
% \Picardfunctor_{X/B} denotes the Picard functor of X over B
\def\Pic{\mathop{\mathrm{Pic}}\nolimits}
\def\Picardstack{\mathcal{P}\!\mathit{ic}}
\def\Picardfunctor{\mathrm{Pic}}
\def\Deformationcategory{\mathcal{D}\!\mathit{ef}}
\usepackage{cleveref}
\crefname{originaldefinition}{original definition}{original definitions}
\Crefname{originaldefinition}{Original definition}{Original definitions}
\makeatletter
\newlabel{formal-spaces-definition-adic-morphism}{{0AQ3}{source}{}{sourcepointer-1}{}}
\newlabel{formal-spaces-definition-adic-morphism@cref}{{[originaldefinition][0][]0AQ3}{[1][0][]source}{}{}{}}
\makeatother
\crefname{originallemma}{original lemma}{original lemmas}
\Crefname{originallemma}{Original lemma}{Original lemmas}
\makeatletter
\newlabel{formal-spaces-lemma-base-change-separated}{{0ARP}{source}{}{sourcepointer-2}{}}
\newlabel{formal-spaces-lemma-base-change-separated@cref}{{[originallemma][0][]0ARP}{[1][0][]source}{}{}{}}
\makeatother
\crefname{originallemma}{original lemma}{original lemmas}
\Crefname{originallemma}{Original lemma}{Original lemmas}
\makeatletter
\newlabel{formal-spaces-lemma-map-completions-representable}{{0APV}{source}{}{sourcepointer-3}{}}
\newlabel{formal-spaces-lemma-map-completions-representable@cref}{{[originallemma][0][]0APV}{[1][0][]source}{}{}{}}
\makeatother
\crefname{originallemma}{original lemma}{original lemmas}
\Crefname{originallemma}{Original lemma}{Original lemmas}
\makeatletter
\newlabel{formal-spaces-lemma-map-into-algebraic-space}{{0AQH}{source}{}{sourcepointer-4}{}}
\newlabel{formal-spaces-lemma-map-into-algebraic-space@cref}{{[originallemma][0][]0AQH}{[1][0][]source}{}{}{}}
\makeatother
\crefname{originallemma}{original lemma}{original lemmas}
\Crefname{originallemma}{Original lemma}{Original lemmas}
\makeatletter
\newlabel{formal-spaces-lemma-separated-local}{{0ARS}{source}{}{sourcepointer-5}{}}
\newlabel{formal-spaces-lemma-separated-local@cref}{{[originallemma][0][]0ARS}{[1][0][]source}{}{}{}}
\makeatother
\crefname{originalreference}{original reference}{original references}
\Crefname{originalreference}{Original reference}{Original references}
\makeatletter
\newlabel{lemma-base-change-rig-surjective}{{0AQS}{source}{}{sourcepointer-6}{}}
\newlabel{lemma-base-change-rig-surjective@cref}{{[originalreference][0][]0AQS}{[1][0][]source}{}{}{}}
\makeatother
\crefname{originalreference}{original reference}{original references}
\Crefname{originalreference}{Original reference}{Original references}
\makeatletter
\newlabel{lemma-monomorphism-rig-surjective}{{0AR0}{source}{}{sourcepointer-7}{}}
\newlabel{lemma-monomorphism-rig-surjective@cref}{{[originalreference][0][]0AR0}{[1][0][]source}{}{}{}}
\makeatother
\crefname{originalreference}{original reference}{original references}
\Crefname{originalreference}{Original reference}{Original references}
\makeatletter
\newlabel{lemma-output-quasi-compact}{{0ARU}{source}{}{sourcepointer-8}{}}
\newlabel{lemma-output-quasi-compact@cref}{{[originalreference][0][]0ARU}{[1][0][]source}{}{}{}}
\makeatother
\crefname{originalreference}{original reference}{original references}
\Crefname{originalreference}{Original reference}{Original references}
\makeatletter
\newlabel{lemma-output-quasi-separated}{{0ARV}{source}{}{sourcepointer-9}{}}
\newlabel{lemma-output-quasi-separated@cref}{{[originalreference][0][]0ARV}{[1][0][]source}{}{}{}}
\makeatother
\crefname{originalreference}{original reference}{original references}
\Crefname{originalreference}{Original reference}{Original references}
\makeatletter
\newlabel{lemma-quasi-separated-gives-quasi-separated}{{0GDG}{source}{}{sourcepointer-10}{}}
\newlabel{lemma-quasi-separated-gives-quasi-separated@cref}{{[originalreference][0][]0GDG}{[1][0][]source}{}{}{}}
\makeatother
\crefname{originalreference}{original reference}{original references}
\Crefname{originalreference}{Original reference}{Original references}
\makeatletter
\newlabel{lemma-separated-mono-open-diagonal-rig-surjective}{{0GDI}{source}{}{sourcepointer-11}{}}
\newlabel{lemma-separated-mono-open-diagonal-rig-surjective@cref}{{[originalreference][0][]0GDI}{[1][0][]source}{}{}{}}
\makeatother
\crefname{originallemma}{original lemma}{original lemmas}
\Crefname{originallemma}{Original lemma}{Original lemmas}
\makeatletter
\newlabel{spaces-cohomology-lemma-check-separated-dvr}{{0ARJ}{source}{}{sourcepointer-12}{}}
\newlabel{spaces-cohomology-lemma-check-separated-dvr@cref}{{[originallemma][0][]0ARJ}{[1][0][]source}{}{}{}}
\makeatother
\crefname{originalremark}{original remark}{original remarks}
\Crefname{originalremark}{Original remark}{Original remarks}
\makeatletter
\newlabel{spaces-cohomology-remark-variant}{{0ARL}{source}{}{sourcepointer-13}{}}
\newlabel{spaces-cohomology-remark-variant@cref}{{[originalremark][0][]0ARL}{[1][0][]source}{}{}{}}
\makeatother
\crefname{originallemma}{original lemma}{original lemmas}
\Crefname{originallemma}{Original lemma}{Original lemmas}
\makeatletter
\newlabel{spaces-morphisms-lemma-base-change-separated}{{03KL}{source}{}{sourcepointer-14}{}}
\newlabel{spaces-morphisms-lemma-base-change-separated@cref}{{[originallemma][0][]03KL}{[1][0][]source}{}{}{}}
\makeatother
\crefname{originalreference}{original reference}{original references}
\Crefname{originalreference}{Original reference}{Original references}
\makeatletter
\newlabel{theorem-dilatations-general}{{0ARB}{source}{}{sourcepointer-15}{}}
\newlabel{theorem-dilatations-general@cref}{{[originalreference][0][]0ARB}{[1][0][]source}{}{}{}}
\makeatother
\title{Separatedness of the algebraization is equivalent to separatedness of the formal model with rig-surjective diagonal}
\author{The Stacks Project Authors}
\date{stacks-project:restricted.tex:local-snapshot-20261012; GNU FDL 1.2 or later}
\begin{document}
\maketitle
\noindent\textit{Editorial note.} All delivered mathematical text is the author's own, copied from the pinned Stacks Project chapter file restricted.tex; no statement or proof step was rewritten, reordered or removed.\par\medskip
\noindent\textit{Editorial note.} The counted proof is the author's complete argument: a forward implication from the chapter's lemmas, an etale-descent reduction to affines, and a three-case valuative criterion closed by lifting along the rig-surjective diagonal and algebraizing, all inside the delivered span\par\medskip
\noindent\textit{Editorial note.} The label lemma-output-separated names the counted statement, and the proof environment immediately following it in the source is delivered verbatim as its proof.\par\medskip
\noindent\textit{Editorial note.} The counted statement is carried by the author's own lemma environment in the chapter source.\par\medskip
\noindent\textit{Allowed standard prerequisites (editorial).}\begin{itemize}
\item The algebraization theorem setting: f : X' to X corresponds to g : W to X\_\{/T\}, with notation as in the chapter's general dilatation theorem
\item The valuative criterion for separatedness of algebraic spaces over a complete DVR, and its formal variant used for adic morphisms of formal algebraic spaces
\item Locality on the base for separatedness under etale base change, for morphisms of spaces and of formal algebraic spaces, and preservation of rig-surjectivity under base change
\end{itemize}\medskip
\ifdefined\hypertarget\hypertarget{sourcepointer-1}{}\fi
\noindent\textit{Source reference.} Original definition 0AQ3: sources/stacks-chow/tags.txt, source lines 13496--13496. Prerequisite pointer invoked at proof line(s) 7268; it names a standard fact of a sibling chapter that the counted argument uses, and every step of the counted proof built on it is retained verbatim in the delivered body.\par\medskip
\ifdefined\hypertarget\hypertarget{sourcepointer-2}{}\fi
\noindent\textit{Source reference.} Original lemma 0ARP: sources/stacks-chow/tags.txt, source lines 13552--13552. Prerequisite pointer invoked at proof line(s) 7217; it names a standard fact of a sibling chapter that the counted argument uses, and every step of the counted proof built on it is retained verbatim in the delivered body.\par\medskip
\ifdefined\hypertarget\hypertarget{sourcepointer-3}{}\fi
\noindent\textit{Source reference.} Original lemma 0APV: sources/stacks-chow/tags.txt, source lines 13488--13488. Prerequisite pointer invoked at proof line(s) 7267; it names a standard fact of a sibling chapter that the counted argument uses, and every step of the counted proof built on it is retained verbatim in the delivered body.\par\medskip
\ifdefined\hypertarget\hypertarget{sourcepointer-4}{}\fi
\noindent\textit{Source reference.} Original lemma 0AQH: sources/stacks-chow/tags.txt, source lines 13510--13510. Prerequisite pointer invoked at proof line(s) 7274; it names a standard fact of a sibling chapter that the counted argument uses, and every step of the counted proof built on it is retained verbatim in the delivered body.\par\medskip
\ifdefined\hypertarget\hypertarget{sourcepointer-5}{}\fi
\noindent\textit{Source reference.} Original lemma 0ARS: sources/stacks-chow/tags.txt, source lines 13555--13555. Prerequisite pointer invoked at proof line(s) 7252; it names a standard fact of a sibling chapter that the counted argument uses, and every step of the counted proof built on it is retained verbatim in the delivered body.\par\medskip
\ifdefined\hypertarget\hypertarget{sourcepointer-6}{}\fi
\noindent\textit{Source reference.} Original reference 0AQS: sources/stacks-restricted/pinned-original-source.tex, source lines 5324--5324. Prerequisite pointer invoked at proof line(s) 7220; it names a standard fact of the same chapter that the counted argument uses, and every step of the counted proof built on it is retained verbatim in the delivered body.\par\medskip
\ifdefined\hypertarget\hypertarget{sourcepointer-7}{}\fi
\noindent\textit{Source reference.} Original reference 0AR0: sources/stacks-restricted/pinned-original-source.tex, source lines 5525--5525. Prerequisite pointer invoked at proof line(s) 7272; it names a standard fact of the same chapter that the counted argument uses, and every step of the counted proof built on it is retained verbatim in the delivered body.\par\medskip
\ifdefined\hypertarget\hypertarget{sourcepointer-8}{}\fi
\noindent\textit{Source reference.} Original reference 0ARU: sources/stacks-restricted/pinned-original-source.tex, source lines 7117--7117. Prerequisite pointer invoked at proof line(s) 7212; it names a standard fact of the same chapter that the counted argument uses, and every step of the counted proof built on it is retained verbatim in the delivered body.\par\medskip
\ifdefined\hypertarget\hypertarget{sourcepointer-9}{}\fi
\noindent\textit{Source reference.} Original reference 0ARV: sources/stacks-restricted/pinned-original-source.tex, source lines 7155--7155. Prerequisite pointer invoked at proof line(s) 7224; it names a standard fact of the same chapter that the counted argument uses, and every step of the counted proof built on it is retained verbatim in the delivered body.\par\medskip
\ifdefined\hypertarget\hypertarget{sourcepointer-10}{}\fi
\noindent\textit{Source reference.} Original reference 0GDG: sources/stacks-restricted/pinned-original-source.tex, source lines 6106--6106. Prerequisite pointer invoked at proof line(s) 7208; it names a standard fact of the same chapter that the counted argument uses, and every step of the counted proof built on it is retained verbatim in the delivered body.\par\medskip
\ifdefined\hypertarget\hypertarget{sourcepointer-11}{}\fi
\noindent\textit{Source reference.} Original reference 0GDI: sources/stacks-restricted/pinned-original-source.tex, source lines 6241--6241. Prerequisite pointer invoked at proof line(s) 7209; it names a standard fact of the same chapter that the counted argument uses, and every step of the counted proof built on it is retained verbatim in the delivered body.\par\medskip
\ifdefined\hypertarget\hypertarget{sourcepointer-12}{}\fi
\noindent\textit{Source reference.} Original lemma 0ARJ: sources/stacks-chow/tags.txt, source lines 13547--13547. Prerequisite pointer invoked at proof line(s) 7228; it names a standard fact of a sibling chapter that the counted argument uses, and every step of the counted proof built on it is retained verbatim in the delivered body.\par\medskip
\ifdefined\hypertarget\hypertarget{sourcepointer-13}{}\fi
\noindent\textit{Source reference.} Original remark 0ARL: sources/stacks-chow/tags.txt, source lines 13549--13549. Prerequisite pointer invoked at proof line(s) 7229; it names a standard fact of a sibling chapter that the counted argument uses, and every step of the counted proof built on it is retained verbatim in the delivered body.\par\medskip
\ifdefined\hypertarget\hypertarget{sourcepointer-14}{}\fi
\noindent\textit{Source reference.} Original lemma 03KL: sources/stacks-chow/tags.txt, source lines 4681--4681. Prerequisite pointer invoked at proof line(s) 7215; it names a standard fact of a sibling chapter that the counted argument uses, and every step of the counted proof built on it is retained verbatim in the delivered body.\par\medskip
\ifdefined\hypertarget\hypertarget{sourcepointer-15}{}\fi
\noindent\textit{Source reference.} Original reference 0ARB: sources/stacks-restricted/pinned-original-source.tex, source lines 7015--7015. Prerequisite pointer invoked at proof line(s) 7199; it names a standard fact of the same chapter that the counted argument uses, and every step of the counted proof built on it is retained verbatim in the delivered body.\par\medskip
\begin{lemma}
\label{lemma-output-separated}
With assumptions and notation as in Theorem \ref{theorem-dilatations-general}
let $f : X' \to X$ correspond to $g : W \to X_{/T}$.
Then $f$ is separated $\Leftrightarrow$ $g$ is
separated and $\Delta_g : W \to W \times_{X_{/T}} W$ is rig-surjective.
\end{lemma}

\begin{proof}
If $f$ is separated, then $g$ is separated and $\Delta_g$
is rig-surjective by
Lemmas \ref{lemma-quasi-separated-gives-quasi-separated} and
\ref{lemma-separated-mono-open-diagonal-rig-surjective}.
Assume $g$ is separated and $\Delta_g$ is rig-surjective.
Exactly as in the proof of
Lemma \ref{lemma-output-quasi-compact}
we may check this over the members of a \'etale covering of $X$
by affine schemes using
Morphisms of Spaces, Lemma \ref{spaces-morphisms-lemma-base-change-separated}
(locality on the base of being separated for morphisms of algebraic spaces),
Formal Spaces, Lemma \ref{formal-spaces-lemma-base-change-separated}
(being separated for morphisms of formal algebraic spaces is preserved
by base change), and
Lemma \ref{lemma-base-change-rig-surjective} (being rig-surjective
is preserved by base change).
Thus we may and do assume $X$ is affine. Furthermore,
we already know that $f : X' \to X$ is quasi-separated by
Lemma \ref{lemma-output-quasi-separated}.

\medskip\noindent
By Cohomology of Spaces, Lemma
\ref{spaces-cohomology-lemma-check-separated-dvr} and
Remark \ref{spaces-cohomology-remark-variant}
it suffices to show that given any commutative diagram
$$
\xymatrix{
\Spec(K) \ar[r] \ar[d] & X' \ar[d] \\
\Spec(R) \ar[r]^p \ar@{-->}[ru] & X' \times_X X'
}
$$
where $R$ is a complete discrete valuation ring with fraction field $K$,
there is a dotted arrow making the diagram commute (as this will
give the uniqueness part of the valuative criterion). Let
$h : \Spec(R) \to X$ be the composition of $p$ with the morphism
$Y \times_X Y \to X$. There are three cases:
Case I: $h(\Spec(R)) \subset U$. This case is trivial
because $U' = X' \times_X U \to U$ is an isomorphism.
Case II: $h$ maps $\Spec(R)$ into $T$. This case follows
from our assumption that $g : W \to X_{/T}$ is separated. Namely,
if $Z$ denotes the reduced induced closed subspace structure
on $T$, then $h$ factors through $Z$ and
$$
W \times_{X_{/T}} Z = X' \times_X Z \longrightarrow Z
$$
is separated by assumption (and for example
Formal Spaces, Lemma \ref{formal-spaces-lemma-separated-local})
which implies we get the lifting property by
Cohomology of Spaces, Lemma \ref{spaces-cohomology-lemma-check-separated-dvr}
applied to the displayed arrow. Case III: $h(\Spec(K))$ is not in $T$
but $h$ maps the closed point of $\Spec(R)$ into $T$. In this case
the corresponding morphism
$$
p_{/T} :
\text{Spf}(R)
\longrightarrow
(X' \times_X X')_{/T} =
W \times_{X_{/T}} W
$$
is an adic morphism (by
Formal Spaces, Lemma
\ref{formal-spaces-lemma-map-completions-representable} and
Definition \ref{formal-spaces-definition-adic-morphism}).
Hence our assumption that
$\Delta_g : W \to W \times_{X_{/T}} W$ is rig-surjective implies we can lift
$p_{/T}$ to a morphism $\text{Spf}(R) \to W = X'_{/T}$, see
Lemma \ref{lemma-monomorphism-rig-surjective}.
Algebraizing the composition $\text{Spf}(R) \to X'$ using
Formal Spaces, Lemma \ref{formal-spaces-lemma-map-into-algebraic-space}
we find a morphism $\Spec(R) \to X'$ lifting $p$ as desired.
\end{proof}

\end{document}
